\section{Simulation Agents: del razonamiento sobre tareas al modelado de procesos sociales}

\subsection{Por qué hacer simulation}
El ponente pasa después a los simulation agents: no solo hacer que el agente resuelva tareas, sino que también interactúe de manera continua en un contexto social, para estudiar el comportamiento colectivo, la difusión de información y los mecanismos de colaboración.

\begin{importantbox}{El valor de simulation no está en «parecerse a los humanos», sino en «poder formar hipótesis verificables»}
Cuando el agente participa como individuo programable en un sistema multiagente, podemos cambiar de manera sistemática la estrategia de memoria, las reglas de comunicación y la función de incentivos, y observar cómo cambia el comportamiento macroscópico.
\end{importantbox}

\subsection{Los tres niveles de mecanismo de la sociedad generativa}
\begin{knowledgebox}{El puente de lo microscópico a lo macroscópico}
\begin{enumerate}
    \item Nivel de memoria: cómo el individuo conserva y recupera sus experiencias;
    \item Nivel de decisión: cómo el individuo actúa dentro de un contexto;
    \item Nivel de interacción: cómo los individuos se influyen entre sí y forman una dinámica colectiva.
\end{enumerate}
Solo cuando estos tres niveles operan de manera conjunta se constituye un «sistema de simulación social susceptible de investigarse».
\end{knowledgebox}

\begin{warningbox}{Riesgo de validez externa de los resultados de la simulación}
Las reglas en un sistema de simulación suelen ser fijadas por quien investiga, de manera que los resultados pueden reflejar más la «elección de reglas» que un mecanismo social real.
Sin una validación suficiente frente a la realidad, las conclusiones de simulation deben considerarse generación de hipótesis, y no demostración de hechos.
\end{warningbox}

\subsection{Resumen de la sección}
Los simulation agents extienden la investigación de Agent al «nivel del comportamiento del sistema». Su producto más valioso no son las historias, sino hipótesis de mecanismo que pueden someterse a prueba experimental y ser refutadas.

